\section{A Formal Model of Bias, Spike Variance, and Tail Variability}
\label{sec:formal_model}

Section~\ref{sec:empirical} established three empirical facts:
(i) activations have a statistically significant mean ($R_x\gg 1$) while gradient
outputs are mean-free ($R_\delta\approx 1$),
(ii) extreme activation entries are dominated by the mean component,
and (iii) mean removal reduces but does not eliminate the spectral spike.
In this section we formalize these observations into a minimal model that
separates \emph{bias}, \emph{structured low-rank variance}, and a
\emph{high-dimensional tail}.
The model is intentionally simple: its purpose is to isolate the mechanisms
needed to explain the empirical phenomena and motivate later algorithmic
interventions.

\subsection{Token-wise representation model}
\label{subsec:model_token}

Let $x_i \in \mathbb{R}^H$ denote the representation (activation) of token $i$
at a fixed layer, and let
\[
X = \begin{bmatrix} x_1^\top \\ \vdots \\ x_N^\top \end{bmatrix}\in\mathbb{R}^{N\times H}
\]
be the stacked activation matrix for a batch of $N$ tokens.
We decompose each token representation as
\begin{equation}
\label{eq:x_decompose}
x_i = \mu + s_i,
\qquad
\mu \in \mathbb{R}^H,
\qquad
\mathbb{E}[s_i]=0,
\end{equation}
where $\mu$ is a shared bias component and $s_i$ captures token-dependent
variation.
Equivalently,
\[
X = M + S,
\qquad
M \triangleq \mathbf{1}\mu^\top,
\qquad
S \triangleq \begin{bmatrix} s_1^\top \\ \vdots \\ s_N^\top \end{bmatrix}.
\]
By construction, mean-centering removes the bias:
\[
X_c \triangleq X - M = S.
\]

\subsection{Spike--tail variance decomposition}
\label{subsec:model_spike_tail}

The empirical results indicate that $S$ (the centered activations) exhibits a
spike--tail spectrum.
We model this by assuming that the centered variation admits an orthogonal
decomposition into a low-rank structured component and a high-dimensional tail:
\begin{equation}
\label{eq:s_decompose}
s_i = s_i^{\text{spike}} + s_i^{\text{tail}},
\qquad
s_i^{\text{spike}} \in \mathcal{U},
\qquad
s_i^{\text{tail}} \in \mathcal{U}^\perp,
\end{equation}
for a low-dimensional subspace $\mathcal{U}\subset\mathbb{R}^H$ with
$\dim(\mathcal{U})=r \ll H$.
Let $P_{\mathcal{U}}$ and $P_{\mathcal{U}^\perp}$ denote the corresponding
orthogonal projectors. Then
\[
s_i^{\text{spike}} = P_{\mathcal{U}}s_i,
\qquad
s_i^{\text{tail}}  = P_{\mathcal{U}^\perp}s_i.
\]

In matrix form,
\[
S = S_{\text{spike}} + S_{\text{tail}},
\qquad
S_{\text{spike}} = SP_{\mathcal{U}},
\qquad
S_{\text{tail}} = SP_{\mathcal{U}^\perp}.
\]

\paragraph{Empirical correspondence.}
In practice, $\mathcal{U}$ is instantiated as the span of the top-$k$ right singular
vectors of the centered activation matrix $X_c$, where $k$ is chosen as a small
fraction of $H$ (e.g., $k=1\%$).
Section~\ref{subsec:spectral_role_bias} showed that mean removal reduces the
leading singular values but leaves a nontrivial residual low-rank variance
subspace, consistent with~\eqref{eq:s_decompose}.

\subsection{Elementwise outliers and the role of bias}
\label{subsec:model_outliers}

The model~\eqref{eq:x_decompose}--\eqref{eq:s_decompose} distinguishes between
two qualitatively different sources of large elementwise magnitudes:
\begin{itemize}
\item \textbf{Bias-driven outliers.} The bias term $\mu$ is dense and shared,
so the mean matrix $M=\mathbf{1}\mu^\top$ can contribute coherently to the
magnitude of many entries.
\item \textbf{Variance-driven extremes.} The random components
$s_i^{\text{spike}}$ and $s_i^{\text{tail}}$ fluctuate across tokens and can
produce large entries through variance, but do not contribute to the mean.
\end{itemize}
Section~\ref{subsec:outlier_mean_dominance} established empirically that, among
the largest $0.1\%$ entries of $X$, squared-magnitude contributions are dominated
by $M$, with smaller contributions from $S_{\text{spike}}$ and negligible
contributions from $S_{\text{tail}}$.

\subsection{Gradient decomposition under mean-biased activations}
\label{subsec:model_grad}

Let $\delta_i \in \mathbb{R}^{H'}$ denote the gradient output (backpropagated
signal) for token $i$ at a linear layer with weight matrix
$W\in\mathbb{R}^{H'\times H}$.
The per-token gradient is
\[
g_i = \delta_i x_i^\top \in \mathbb{R}^{H'\times H},
\]
and the batch gradient is
\[
G \triangleq \frac{1}{N}\sum_{i=1}^N g_i
= \mathbb{E}[\delta x^\top].
\]
Substituting~\eqref{eq:x_decompose} yields the identity
\begin{equation}
\label{eq:G_decompose}
G
=
\mathbb{E}[\delta]\mu^\top
+
\mathbb{E}[\delta s^\top].
\end{equation}
This separates a rank-one bias-driven term from a cross-covariance term.
Section~\ref{subsec:first_order_asymmetry} established empirically that
$\mathbb{E}[\delta]$ is noise-level, while $\mu$ is large.
Section~\ref{subsec:grad_spectrum_bias} further showed that centering $x$
reduces the spectral spike of $G$, consistent with removing the first term.

\subsection{Second-moment (Adam) statistics}
\label{subsec:model_adam}

Adaptive optimizers maintain elementwise second moments of gradients.
For the outer-product gradient $g_i=\delta_i x_i^\top$, the elementwise second
moment admits the exact batch identity
\begin{equation}
\label{eq:v_identity}
v \triangleq \mathbb{E}[g_i \odot g_i]
=
\frac{1}{N}(\Delta^{\odot 2})^\top (X^{\odot 2}),
\end{equation}
where $\Delta\in\mathbb{R}^{N\times H'}$ stacks $\delta_i$ and
$X\in\mathbb{R}^{N\times H}$ stacks $x_i$.
Equation~\eqref{eq:v_identity} makes explicit that activation magnitudes
directly influence second-moment inflation.
Section~\ref{subsec:optimizer_bias} showed empirically that mean removal reduces
the extreme tail of $v$ (up to $50\times$ in $\max(v)$ in some regimes),
establishing bias-driven second-moment inflation as a key mechanism.

\paragraph{Summary.}
The model isolates three components with distinct roles:
\begin{enumerate}
\item $\mu$ (bias): drives elementwise outliers and inflates second moments;
\item $\mathcal{U}$ (spike variance subspace): captures low-rank structured variance;
\item $\mathcal{U}^\perp$ (tail): carries high-dimensional token-dependent variability.
\end{enumerate}
The next section studies algorithmic consequences: mean subtraction removes
$\mu$ exactly with negligible cost, while lightweight spectral conditioning can
reduce residual spike variance when needed.